\documentclass[UTF8,a4paper,10pt,fontset=none]{ctexart}
\usepackage[margin=18mm,headheight=15pt]{geometry}
\usepackage{amsmath,booktabs,array,tabularx,xcolor,fancyhdr,hyperref}
\setmainfont{Arial}
\setCJKmainfont{Microsoft YaHei}
\setCJKmonofont{Microsoft YaHei}
\definecolor{ink}{HTML}{18344C}
\hypersetup{unicode=true,colorlinks=true,urlcolor=ink,linkcolor=ink,pdftitle={YunTAPR-Net V1.3 极端降水审计结果汇总},pdfauthor={YunTAPR-Net Research}}
\pagestyle{fancy}\fancyhf{}\lhead{\small YunTAPR-Net V1.3}\rhead{\small 极端降水专项审计}\cfoot{\small\thepage}
\setlength{\parindent}{0pt}\setlength{\parskip}{4pt}
\setlength{\tabcolsep}{4pt}\renewcommand{\arraystretch}{1.13}
\ctexset{section={format=\large\bfseries,beforeskip=8pt,afterskip=4pt}}
\emergencystretch=2em
\newcommand{\result}[1]{\textcolor{ink}{\textbf{#1}}}
\begin{document}
{\LARGE\bfseries YunTAPR-Net V1.3}\par
{\Large 极端降水审计结果汇总}\par
2026年10月31日\quad 2024 Development\quad 研究者主导，AI辅助证据整理

\section*{1. 先看结论}
\result{当前没有证据表明B7/B8已经取得稳定的强雨点检出增益。}
九个模型配置在A/B/C三种点估计下，严格$R>20$ mm/h均为
\[
 TP=0,\quad FP=0,\quad FN=717,\quad CSI=POD=0.
\]
这些717个像元时次来自相同标签人口，不能乘以模型或seed数量当成新增独立样本。

\textbf{有限正面结果：}B7相对B4在20--30和30--50 mm/h的条件Pinball各有2/3 seed改善；208个强雨发生窗口的总体Brier改善3/3，但窗口还含弱雨和无雨像元，不能等同于极端雨强校准改善。

\textbf{主要反例：}B8相对B7在上述两个强雨层均仅1/3 seed改善，平均Pinball退化；A方案5 mm/h的CSI也仅1/3 seed改善。总体AUROC收益不能替代极端检出、尾部校准或物理机制证据。

\section*{2. 样本、身份与单位}
\begin{center}
\small
\begin{tabular}{@{}lp{116mm}@{}}
\toprule
项目 & 已核实结果 \\
\midrule
模型/随机性 & B4、B7、B8；2029/2030/2031；共9个模型配置 \\
训练与评价 & E0、固定E9 LAST；2023 Train=959、2024 Dev=963 \\
有效监督/有雨像元时次 & 3,303,090 / 407,820（有雨定义$R>0.1$） \\
$R>5$ / $R>10$ & 25,661 / 5,516个像元时次 \\
$R>20$ & 717个像元时次、207个采样窗口 \\
$R\ge20$ & 至少718个像元时次、208个采样窗口 \\
$R>50$ & 6个像元时次、5个采样窗口 \\
独立极端天气事件数 & 不可识别，不能用像元/窗口/周块数量代替 \\
\bottomrule
\end{tabular}
\end{center}

九份E9权重SHA256均与索引及原收据一致；共同人口、掩膜、有效性口径与同seed的E9训练顺序一致。使用GADM4.1中心点行政区mask，SP04裁剪100$\times$100、3430有效点；GFS固定f006、回顾估计，分析时刻00/06/12/18 UTC。

\result{阈值边界不能混用：}$>20$有207个窗口，$\ge20$有208个，至少一个窗口最大值恰为20。缺少标签/预测等于20、50的联合计数，不能把原来的$(20,50]$改称$[20,50)$。

IMERG保存单位为mm/h，30分钟累计量=雨强$\times0.5$ h；5/10/20/50 mm/h分别对应2.5/5/10/25 mm每30分钟，不能再把现有雨强标签乘2。

本次只读取已有证据；新增训练、模型推理和2025载荷访问均为0，原始数据及既有结果未覆盖。
\newpage
\section*{3. 超阈检出：A方案逐seed结果}
A=$p\,\mathrm{mean}(\exp(q)-1)$为诊断代理，\textbf{未科学冻结}。下表严格比较Pred$>r$与Truth$>r$；CSI/POD/FAR/Precision均以百分比显示。B/C原值及所有同seed差值保留在PDF附件CSV。
\begin{center}
\footnotesize
\begin{tabular}{@{}lrrrrrrrrr@{}}
\toprule
模型 & seed & $r$ & TP & FP & FN & CSI(\%) & POD(\%) & FAR(\%) & Precision(\%) \\
\midrule
B4 & 2029 & 5 & 718 & 268 & 24943 & 2.7691 & 2.7980 & 27.1805 & 72.8195 \\
B4 & 2029 & 10 & 1 & 2 & 5515 & 0.0181 & 0.0181 & 66.6667 & 33.3333 \\
B4 & 2029 & 20 & 0 & 0 & 717 & 0.0000 & 0.0000 & 未定义 & 未定义 \\
B4 & 2030 & 5 & 19 & 20 & 25642 & 0.0740 & 0.0740 & 51.2821 & 48.7179 \\
B4 & 2030 & 10 & 0 & 0 & 5516 & 0.0000 & 0.0000 & 未定义 & 未定义 \\
B4 & 2030 & 20 & 0 & 0 & 717 & 0.0000 & 0.0000 & 未定义 & 未定义 \\
B4 & 2031 & 5 & 267 & 158 & 25394 & 1.0341 & 1.0405 & 37.1765 & 62.8235 \\
B4 & 2031 & 10 & 0 & 0 & 5516 & 0.0000 & 0.0000 & 未定义 & 未定义 \\
B4 & 2031 & 20 & 0 & 0 & 717 & 0.0000 & 0.0000 & 未定义 & 未定义 \\
B7 & 2029 & 5 & 249 & 88 & 25412 & 0.9670 & 0.9703 & 26.1128 & 73.8872 \\
B7 & 2029 & 10 & 0 & 1 & 5516 & 0.0000 & 0.0000 & 100.0000 & 0.0000 \\
B7 & 2029 & 20 & 0 & 0 & 717 & 0.0000 & 0.0000 & 未定义 & 未定义 \\
B7 & 2030 & 5 & 139 & 228 & 25522 & 0.5369 & 0.5417 & 62.1253 & 37.8747 \\
B7 & 2030 & 10 & 0 & 4 & 5516 & 0.0000 & 0.0000 & 100.0000 & 0.0000 \\
B7 & 2030 & 20 & 0 & 0 & 717 & 0.0000 & 0.0000 & 未定义 & 未定义 \\
B7 & 2031 & 5 & 392 & 476 & 25269 & 1.4998 & 1.5276 & 54.8387 & 45.1613 \\
B7 & 2031 & 10 & 1 & 21 & 5515 & 0.0181 & 0.0181 & 95.4545 & 4.5455 \\
B7 & 2031 & 20 & 0 & 0 & 717 & 0.0000 & 0.0000 & 未定义 & 未定义 \\
B8 & 2029 & 5 & 316 & 115 & 25345 & 1.2259 & 1.2314 & 26.6821 & 73.3179 \\
B8 & 2029 & 10 & 0 & 1 & 5516 & 0.0000 & 0.0000 & 100.0000 & 0.0000 \\
B8 & 2029 & 20 & 0 & 0 & 717 & 0.0000 & 0.0000 & 未定义 & 未定义 \\
B8 & 2030 & 5 & 136 & 200 & 25525 & 0.5259 & 0.5300 & 59.5238 & 40.4762 \\
B8 & 2030 & 10 & 0 & 6 & 5516 & 0.0000 & 0.0000 & 100.0000 & 0.0000 \\
B8 & 2030 & 20 & 0 & 0 & 717 & 0.0000 & 0.0000 & 未定义 & 未定义 \\
B8 & 2031 & 5 & 104 & 193 & 25557 & 0.4023 & 0.4053 & 64.9832 & 35.0168 \\
B8 & 2031 & 10 & 0 & 4 & 5516 & 0.0000 & 0.0000 & 100.0000 & 0.0000 \\
B8 & 2031 & 20 & 0 & 0 & 717 & 0.0000 & 0.0000 & 未定义 & 未定义 \\
\bottomrule
\end{tabular}
\end{center}
FAR=$FP/(TP+FP)$，Precision=$TP/(TP+FP)$。没有预测阳性时二者\textbf{未定义}，不能填0或解释为零虚警优势。$>50$ mm/h也可由现有计数严格补出：27组配置/点方案均TP=FP=0、FN=6、TN=3,303,084。

\textbf{A方案CSI的配对变化：}实验组减参考组；pp为百分点，正值较好。
\begin{center}
\small
\begin{tabular}{@{}lrrrr@{}}
\toprule
比较 & $r$ (mm/h) & $\Delta$CSI均值(pp) & 样本SD(pp) & 改善seed \\
\midrule
B7$-$B4 & 5 & -0.29116 & 1.30849 & 2/3 \\
B7$-$B4 & 10 & -0.00002 & 0.01809 & 1/3 \\
B8$-$B4 & 5 & -0.57437 & 0.99877 & 1/3 \\
B8$-$B4 & 10 & -0.00604 & 0.01046 & 0/3 \\
B8$-$B7 & 5 & -0.28321 & 0.71802 & 1/3 \\
B8$-$B7 & 10 & -0.00602 & 0.01043 & 0/3 \\
\bottomrule
\end{tabular}
\end{center}
$>20$的CSI差值全部为0；这表示共同没有命中，不表示模型达到相同且可接受的极端技能。固定概率阈值的判别与点估计检出也不能互换。
\newpage
\section*{4. 20--50 mm/h分位数误差与反例}
Pinball为有雨条件、冻结log1p尺度损失，越低越好，\textbf{不是mm/h的MAE}。七层原值均在附件；这里突出20--50及稀少的$>50$层。
\begin{center}
\footnotesize
\begin{tabular}{@{}lrrrrr@{}}
\toprule
模型 & seed & Pinball(20,30] & Pinball(30,50] & Pinball$>50$ & q32(20,50](\%) \\
\midrule
B4 & 2029 & 0.832683 & 0.986955 & 1.350198 & 缺明细 \\
B4 & 2030 & 0.839462 & 1.010071 & 1.325037 & 0.703 \\
B4 & 2031 & 0.818505 & 0.974025 & 1.313961 & 2.813 \\
B7 & 2029 & 0.861780 & 1.019434 & 1.357232 & 5.485 \\
B7 & 2030 & 0.808730 & 0.970630 & 1.289152 & 2.532 \\
B7 & 2031 & 0.796965 & 0.958989 & 1.271548 & 7.173 \\
B8 & 2029 & 0.832745 & 0.988714 & 1.337625 & 5.767 \\
B8 & 2030 & 0.812781 & 0.974595 & 1.317422 & 1.969 \\
B8 & 2031 & 0.835951 & 0.997748 & 1.295502 & 1.266 \\
\bottomrule
\end{tabular}
\end{center}
$(20,30]$、$(30,50]$、$>50$分别有564、147、6个像元时次。q32列针对$(20,50]$共711个像元：按真实雨强条件化的覆盖仅作描述，不要求每个雨强层达到98.4375\%名义水平。B4/2029缺逐场景尾部明细，不能填补其缺失覆盖率。

\textbf{同seed差值汇总：}负值改善；SD为三个配对差值的样本标准差。
\begin{center}
\small
\begin{tabular}{@{}lrrrr@{}}
\toprule
比较 & 雨强层(mm/h) & $\Delta$Pinball均值 & 样本SD & 改善seed \\
\midrule
B7$-$B4 & (20,30] & -0.00772516 & 0.03221849 & 2/3 \\
B7$-$B4 & (30,50] & -0.00733308 & 0.03657387 & 2/3 \\
B8$-$B4 & (20,30] & -0.00305797 & 0.02222877 & 1/3 \\
B8$-$B4 & (30,50] & -0.00333161 & 0.02992591 & 1/3 \\
B8$-$B7 & (20,30] & +0.00466719 & 0.03401516 & 1/3 \\
B8$-$B7 & (30,50] & +0.00400148 & 0.03473996 & 1/3 \\
\bottomrule
\end{tabular}
\end{center}
改善次数、均值和标准差应一起看，不能只选有利seed或层。三seed是训练随机性的描述性复现，不能证明统计显著、跨年份泛化或独立天气事件稳健性。

\section*{5. 峰值低估：能证明下界，不能伪造精确误差}
QC真值峰值为68.0599976 mm/h。所有27组点估计都没有超过20的float32预测，因此真值峰值位置的低估\textbf{至少约48.06 mm/h}。为覆盖float64转float32舍入，实际计算上界使用nextafter(threshold,$+\infty$)。部分配置的下界更高，见附件。

这只是由保存计数推出的\textbf{下界}，不是精确峰值误差，也不能用下界给模型排名。强雨条件MAE/RMSE/Bias、逐事件预测峰值和峰值位置没有联合明细；已有连续误差只针对全体有效像元，不能冒充强雨误差。
\newpage
\section*{6. 归因、概率识别与仍缺的证据}
\begin{tabularx}{\linewidth}{@{}p{26mm}XX@{}}
\toprule
层次 & 现有证据 & 不能据此声称\\
\midrule
事件 & 963采样窗口，208个窗口最大值$\ge20$；$>50$的6像元在5窗口。 & 没有事件ID/追踪，不能称208或5场独立暴雨。\\
\addlinespace
高程/风/水汽 & B7/B8有三个高程边际；B4对应高程未保存；缺迎/背风和PWAT联合误差。 & 不能从边际总体误差推出高山极端或迎风抬升收益。\\
\addlinespace
DOTE & 真实数值对ZERO三seed的Brier改善3/3，但排序及强雨收益不稳；ZERO有128权重无损失梯度；PERP仅单seed。 & B7减B4还包含静态地形、参数和融合变化；负对照亦有容量/分布/空间混杂，未证物理因果。\\
\addlinespace
DTFM & B8/B7同信息、同名义活跃参数量、同seed，比较的是替换融合形式。 & 总体AUROC改善不证明极端命中、尾部校准或最终年份泛化。\\
\bottomrule
\end{tabularx}

\textbf{极端概率必须用完整混合语义：}
\[
P(R>r)=p_{\rm occ}\,[1-F_{\rm cond}(\log(1+r))].
\]
不能复用$P(R>0.1)\ge0.5$。既有D表的q32支持域因模型不同而变化，原支持子集率不能公平排名。附件另给\texttt{D\_BOUND\_PGE05}：在$q32$及以上，任何与$\tau_{32}=63/64$一致的单调尾部补全均有$P(R>r)\le1/64<0.5$，所以可补全同人口的0.5决策计数；\textbf{这没有恢复精确尾部概率、极端Brier/AUROC/AP}，也没有将未知概率设为0。

\section*{7. 下一步建议：新增训练0，先补评价缺口}
\begin{enumerate}
\item 先冻结独立事件/时间依赖定义，以及$\ge20$、$[20,50)$、$\ge50$的边界；保留原严格$>$结果作对照。
\item 从2023 Train同人口、同掩膜正雨标签确定P95/P99；明确$R>0$或$R>0.1$口径。当前没有可用完整正雨分布，不提供猜测数值，不用Dev或2025选择阈值。
\item 后续如获相应评价预算，只对\textbf{九份既有E9}各一次2024 Dev前向，保存必要的逐像元/逐场景联合量，补强雨连续误差、精确峰值及高程/风/水汽交叉统计；另需独立事件目录，四UTC前向本身不能补足天气系统追踪。
\item 暂不重训B0--B8、不加模块、不选最终模型。任何新增训练或2025揭盲仍需研究者另行明确授权。
\end{enumerate}

\section*{8. 来源与随文件保存的明细}
汇总依据已完成只读审计：109份定位源文件、九份E9身份、1400行CSV。108组点分类、27组混合概率决策界、1438个配对指标单元格已核对；492条seed统计已检验样本标准差口径。此汇总未重算原科研指标。

源交付提交：\href{https://github.com/flyeagle-cex/YunTAPR-Net/commit/cfe525b30db6fc7c36c578d7904da9238d71e7cd}{\texttt{cfe525b}}。PDF附件含\textbf{审计明细.csv、完整审计报告.md、文档源码.tex}；源码可提取后用XeLaTeX编辑编译。附件为脱敏版，原完整私有路径版本继续保存在本地审计交付中。

{\scriptsize CSV SHA256：\texttt{99a2be1aae0aadc725cdc0e9e2c0149e}\par\texttt{0e6ca7b7b3feb6c885382c7cb255c062}\par
审计报告SHA256：\texttt{b001d806d1d5989a3b6e24fd03549c14}\par\texttt{fe14b7999b1db5151d462774ccfee40e}}
\end{document}
